\documentclass[11pt]{article}
\usepackage[margin=1in]{geometry}
\usepackage{amsmath}
\usepackage{cite}
\usepackage{hyperref}

\begin{document}

\begin{flushleft}
\textbf{Name:} Kevin Beltran \\
\textbf{Course:} CSC 570 Computational Geometry \\
\textbf{Assignment:} Project Proposal \\
\textbf{Date:} February 10, 2026
\end{flushleft}

\begin{center}
    \section*{Experimental Comparison of Visvalingam-Whyatt and Douglas-Peucker Line Simplification}
\end{center}

\subsection*{1. Statement of the Proposed Topic}
This project will experimentally evaluate and compare two primary algorithms for 2D polyline simplification: the \textbf{Douglas-Peucker} algorithm and the \textbf{Visvalingam-Whyatt} algorithm. The primary goal is to analyze how different geometric error metrics (perpendicular distance versus triangle area) impact the "visual character" and topological integrity of a curve under high compression. This  is particularly relevant for mobile application development, where GPS tracks must be simplified for efficient storage and real-time rendering on resource-constrained devices without losing critical geospatial features.

\subsection*{2. Brief Survey of Known Results}
\begin{itemize}
    \item \textbf{Douglas-Peucker (1973):} A top-down recursive algorithm that uses perpendicular distance as its threshold. While highly effective at preserving shape, it is prone to "spiking" artifacts and potential self-intersections during extreme simplification.
    \item \textbf{Visvalingam-Whyatt (1992):} A bottom-up "greedy" algorithm that iteratively removes the vertex forming the triangle with the smallest "effective area." It is noted for preserving the "bulk" of organic shapes, such as coastlines and rivers.
    \item \textbf{Geometric Primitives:} Douglas-Peucker relies on the point-to-line distance formula. In contrast, Visvalingam-Whyatt utilizes the shoelace formula or cross-product to calculate triangle area $(p_{i-1}, p_i, p_{i+1})$.
    \item \textbf{Hybrid Approaches:} Industry-standard libraries often utilize a fast $O(n)$ Radial Distance filter as a pre-processing step to reduce the input size before running more computationally expensive $O(n \log n)$ algorithms.
\end{itemize}

\subsection*{3. Potential Plan of Attack}
\begin{enumerate}
    \item \textbf{Implementation:} I will implement the standard recursive Douglas-Peucker and a min-heap-accelerated Visvalingam-Whyatt in Python to ensure $O(n \log n)$ efficiency for both.
    \item \textbf{Data Sourcing:} I will source open-source GeoJSON data representing diverse geometries: San Luis Obispo hiking trails (mechanical/linear) and complex coastal boundaries (organic/fractal).
    \item \textbf{Benchmarking:} I will generate comparative graphs measuring \textit{Compression Ratio} (vertices removed) versus \textit{Geometric Error}. Error will be quantified using both maximum perpendicular distance and \textbf{Hausdorff distance} to measure deviation between the original and simplified sets.
    \item \textbf{Half-Baked Idea:} I'd like to also expore some "Shape-Aware Selector" heuristic. For example, a pre-pass to calculate the local sinuosity (curviness) of segments, automatically switching between Douglas-Peucker for straight segments and Visvalingam-Whyatt for highly curved organic sections to optimize visual fidelity.
\end{enumerate}

\end{document}